\section{The Induction Proof of the Main Theorem: \texorpdfstring{$2n\ge 10$}{2n >= 10}}
\label{sec: 2n >= 10}
\label{sec: high arity induction}

We now complete the induction framework of
\cref{subsec: Induction framework}.  The normalization from
\cref{lm: unitary wire normalization} is valid in every arity at least
eight.  At arity at least ten, however, its strong equality outcome is
combinatorially impossible.

\begin{lemma}
\label{lm: high arity equality rigidity}
No nonzero signature of arity $2n\geq10$ can simultaneously satisfy
{\sc 2nd-Orth} and the strong equality property.
\end{lemma}

\begin{proof}
Similar to \cref{def: f8 matching}, for any pair of distinct indices $e=\{i,j\}$, we can define $M_e, F_e$ and the relation $\sim.$
Also similar to \cref{lm: f8 matching equivalence}, $\sim$ is an equivalence relation.
The proof there makes no special use of $2n=8$, and it holds for any $2n\ge10.$
Consequently the perfect matchings $F_e$ partition $E(K_{2n})$ into a
$1$-factorization.

Choose distinct perfect matchings $A,B$ and disjoint edges $e\in A$, $d\in B$.
Such a choice exists because $n\geq5$.  
Because $A$ and $B$ are distinct, they are disjoint as two equivalent classes of $\sim.$
Compute $\partial_d\partial_eh$ in the two orders.  Starting with $e$, the two
endpoints of $d$ lie in two distinct edges of $A\setminus\{e\}$.
Contracting $d$ switches those two edges into one edge and leaves the
other $m-3$ edges of $A$ unchanged.  Hence the residual matching consists of one switched edge and $m-3$ unchanged $A$-edges.  Reversing the order
gives one switched edge and $m-3$ unchanged $B$-edges.

The two double contractions are the same nonzero equality product, so
their residual matchings must coincide.  
%Distinct factors of a $1$-factorization share no edge. 
Since $m\geq5$, the first residual
matching contains at least two unchanged $A$-edges.  
Neither can be an unchanged $B$-edge, since $A$ and $B$ are disjoint.
However, at most one can equal the single switched edge
in the reverse-order matching.  This is impossible.  Therefore the two
hypotheses cannot hold simultaneously.
\end{proof}

The following descent theorem plays the same role as
\cite[Lemma~9.1]{realholant}. 

\begin{theorem}
\label{thm: strict arity descent}
Let $\mathcal F=\overline{\mathcal F}$ fail condition
\rm{(\ref{cond: tractable classes})}, and let $f\in \mathcal{F}$ is a
signature of arity $2n\geq10$ with
$f\notin\mathcal U^{\otimes}$.  Then either
$\Holant(\mathcal F)$ is \#P-hard, or we can realize an even arity signature
$g\notin\mathcal U^{\otimes}$ of arity at most $2n-2$.
\end{theorem}

\begin{proof}
If $f$ is reducible, \cref{lm: decomposition} makes its proper factors
available.  A nonzero odd-arity factor gives hardness by
\cref{lm: odd holant with conjugation}.  If all factors have even arity,
at least one proper factor lies outside $\mathcal U^{\otimes}$, and it is
the desired $g$.

Assume therefore that $f$ is irreducible.  If it fails {\sc 2nd-Orth},
\cref{lm: not 2nd-orth is hard} gives hardness.  If some equality
contraction $\partial_{ij}f$ lies outside
$\mathcal U^{\otimes}$, take $g=\partial_{ij}f$, whose arity is
$2n-2$.

The only remaining case is that $f$ is irreducible, satisfies
{\sc 2nd-Orth}, and every $\partial_{ij}f$ belongs to
$\mathcal U^{\otimes}$.  Apply
\cref{lm: unitary wire normalization}.  Its first outcome is precisely the
required lower-arity $g$.  Its second outcome is a nonzero arity-$2n$
signature satisfying both {\sc 2nd-Orth} and the strong equality property,
contradicting \cref{lm: high arity equality rigidity}.  Thus one of the
first two outcomes must occur.
\end{proof}

The induction below is essentially the same as the final induction in
\cite[Theorem~9.2]{realholant}, with the complex arities $6$ and $8$
supplied by \cref{lm: f6 is hard,thm: base case 2n=8}.  We include it to
make the proof of the main theorem complete.

\begin{theorem}[Main theorem]
\label{thm: conjugate closed Holant dichotomy}
Let $\mathcal F$ be a conjugate-closed set of algebraic complex-valued Boolean
signatures.  
If $\mathcal F$ satisfies condition
\rm{(\ref{cond: tractable classes})}, then $\Holant(\mathcal F)$ is
polynomial-time computable.  Otherwise $\Holant(\mathcal F)$ is
\#P-hard.
\end{theorem}

\begin{proof}
The tractable direction is \cref{thm: tractable classes}.  Suppose that
$\mathcal F$ does not satisfy condition
\rm{(\ref{cond: tractable classes})}.  If it contains a nonzero odd-arity
signature, hardness follows from
\cref{lm: odd holant with conjugation}.  We may therefore restrict to
nonzero even-arity signatures.

Since $\mathcal U^{\otimes}\subseteq\mathscr T$, the assumption that
$\mathcal F$ fails condition \rm{(\ref{cond: tractable classes})} implies
that some $f\in\mathcal F$ has even arity $2n$ and
$f\notin\mathcal U^{\otimes}$.  We induct on $2n$.  The cases
$2n=2,4,6,8$ are respectively
\cref{lm: base case 2n=2,lm: base case 2n=4,lm: induction 2n=6,%
thm: base case 2n=8}.

For $2n\geq10$, apply \cref{thm: strict arity descent}.  Either hardness
already follows, or it supplies an available
$g\notin\mathcal U^{\otimes}$ of smaller even arity.  
Since $\mathcal{F}$ is conjugate closed, $\overline{g}$ is also realizable.
Let $\mathcal{G}=\mathcal{F}\cup\{g,
\overline{g}\}$.
Clearly $\mathcal{G}$ fails
condition \rm{(\ref{cond: tractable classes})} because it contains $\mathcal{F}.$
The induction hypothesis makes the augmented problem $\Holant(\mathcal{G})$
\#P-hard, while gadget realization reduces it to
$\Holant(\mathcal F)$.  
Since each descent lowers the fixed signature
arity by at least two, the process reaches one of the four base arities
after finitely many constant-size reductions.  This completes the proof.
\end{proof}
